% GENERATED FILE. Edit canonical lessons, then run npm run generate:books.
% canonical-source-hash: fnv1a64:975938f0488a8fca
% canonical-lessons: HI-C314-afvah, HI-C314-carca, HI-C314-dahsat, HI-C314-bharkana, HI-C314-sansani

\chapter{A Rumour Spreads}
\label{ch:hi-a-rumour-spreads}
\hlchaptermodality{314}

\begin{chapteropening}
\textbf{By the end of this chapter:} \emph{I can say how a rumour spread and how people reacted.}

Complete the last lesson of chapter 314: I can say how a rumour spread and how people reacted.
\end{chapteropening}

\section[\texorpdfstring{\hi{अफ़वाह} (afvāh)}{अफ़वाह (afvāh)}]{\hi{अफ़वाह} (afvāh) — a rumour}

\label{lesson:HI-C314-afvah}



\begin{quote}
Before the new one: say the Hindi for an editor, then the Hindi for a detailed account.
\end{quote}

\subsection*{You'll want to know: \hi{अफ़वाह}}
\textbf{\hi{अफ़वाह}} — \emph{afvāh} — \textquotedblleft{}a rumour\textquotedblright{}.

\begin{quote}
\hi{बाज़ार} \hi{में} \hi{एक} \hi{अफ़वाह} \hi{फैली।}
\end{quote}

\begin{quote}
\emph{bāzār meṁ ek afvāh phailī.}
\end{quote}

\begin{quote}
A rumour spread through the market.
\end{quote}

Feminine.

\begin{cousinweb}
\textbf{\hi{अफ़वाह}} (\emph{afvāh}) is usually traced, through Persian, to Arabic \emph{afwāh}, \textquotedblleft{}mouths\textquotedblright{}: what goes from mouth to mouth. It is feminine: \textbf{\hi{अफ़वाह} \hi{फैली}}.
\end{cousinweb}

\begin{grammarlens}[title={What you've built}]
One more word for telling what happened.
\end{grammarlens}

\subsection*{Your turn}
\begin{itemize}
\raggedright
  \item \emph{Say it:} \emph{afvāh}
  \item \emph{Say it:} \emph{afvāh}, once more
  \item \emph{Say it:} \emph{bāzār meṁ ek afvāh phailī.} — A rumour spread through the market.
  \item \emph{Recall:} say the Hindi for an editor, then the Hindi for a detailed account, then say \emph{afvāh} again
\end{itemize}

\begin{tcolorbox}[breakable,colback=teal!4,colframe=teal!35!black,title={Before you move on}]
What does \hi{अफ़वाह} mean? (\textquotedblleft{}A rumour\textquotedblright{}.) Say it once more.
\end{tcolorbox}
\section[\texorpdfstring{\hi{चर्चा} (carcā)}{चर्चा (carcā)}]{\hi{चर्चा} (carcā) — talk; a discussion}

\label{lesson:HI-C314-carca}



\begin{quote}
Before the new one: say the Hindi for a detailed account, then the Hindi for a rumour.
\end{quote}

\subsection*{You'll want to know: \hi{चर्चा}}
\textbf{\hi{चर्चा}} — \emph{carcā} — \textquotedblleft{}talk; a discussion\textquotedblright{}.

\begin{quote}
\hi{पूरे} \hi{शहर} \hi{में} \hi{उस} \hi{अफ़वाह} \hi{की} \hi{चर्चा} \hi{थी।}
\end{quote}

\begin{quote}
\emph{pūre śahar meṁ us afvāh kī carcā thī.}
\end{quote}

\begin{quote}
The whole town was talking about that rumour.
\end{quote}

Feminine.

\begin{cousinweb}
\textbf{\hi{चर्चा}} (\emph{carcā}) is Sanskrit \emph{carcā} (Monier-Williams), a repeating of words and a thinking over; in Hindi it is what everyone is talking about.
\end{cousinweb}

\begin{grammarlens}[title={What you've built}]
One more word for telling what happened.
\end{grammarlens}

\subsection*{Your turn}
\begin{itemize}
\raggedright
  \item \emph{Say it:} \emph{carcā}
  \item \emph{Say it:} \emph{carcā}, once more
  \item \emph{Say it:} \emph{pūre śahar meṁ us afvāh kī carcā thī.} — The whole town was talking about that rumour.
  \item \emph{Recall:} say the Hindi for a detailed account, then the Hindi for a rumour, then say \emph{carcā} again
\end{itemize}

\begin{tcolorbox}[breakable,colback=teal!4,colframe=teal!35!black,title={Before you move on}]
What does \hi{चर्चा} mean? (\textquotedblleft{}Talk; a discussion\textquotedblright{}.) Say it once more.
\end{tcolorbox}
\section[\texorpdfstring{\hi{दहशत} (dahśat)}{दहशत (dahśat)}]{\hi{दहशत} (dahśat) — terror; panic}

\label{lesson:HI-C314-dahsat}



\begin{quote}
Before the new one: say the Hindi for a rumour, then the Hindi for talk.
\end{quote}

\subsection*{You'll want to know: \hi{दहशत}}
\textbf{\hi{दहशत}} — \emph{dahśat} — \textquotedblleft{}terror; panic\textquotedblright{}.

\begin{quote}
\hi{अफ़वाह} \hi{से} \hi{लोगों} \hi{में} \hi{दहशत} \hi{फैल} \hi{गई।}
\end{quote}

\begin{quote}
\emph{afvāh se logoṁ meṁ dahśat phail gaī.}
\end{quote}

\begin{quote}
The rumour spread panic among the people.
\end{quote}

Feminine.

\begin{cousinweb}
\textbf{\hi{दहशत}} (\emph{dahśat}) is usually traced, through Persian, to Arabic \emph{dahsha}, \textquotedblleft{}bewilderment, alarm\textquotedblright{}. It is a fear that spreads through a whole place.
\end{cousinweb}

\begin{grammarlens}[title={What you've built}]
One more word for telling what happened.
\end{grammarlens}

\subsection*{Your turn}
\begin{itemize}
\raggedright
  \item \emph{Say it:} \emph{dahśat}
  \item \emph{Say it:} \emph{dahśat}, once more
  \item \emph{Say it:} \emph{afvāh se logoṁ meṁ dahśat phail gaī.} — The rumour spread panic among the people.
  \item \emph{Recall:} say the Hindi for a rumour, then the Hindi for talk, then say \emph{dahśat} again
\end{itemize}

\begin{tcolorbox}[breakable,colback=teal!4,colframe=teal!35!black,title={Before you move on}]
What does \hi{दहशत} mean? (\textquotedblleft{}Terror; panic\textquotedblright{}.) Say it once more.
\end{tcolorbox}
\section[\texorpdfstring{\hi{भड़काना} (bhaṛkānā)}{भड़काना (bhaṛkānā)}]{\hi{भड़काना} (bhaṛkānā) — to stir up; to incite}

\label{lesson:HI-C314-bharkana}



\begin{quote}
Before the new one: say the Hindi for talk, then the Hindi for terror.
\end{quote}

\subsection*{You'll want to know: \hi{भड़काना}}
\textbf{\hi{भड़काना}} — \emph{bhaṛkānā} — \textquotedblleft{}to stir up; to incite\textquotedblright{}.

\begin{quote}
\hi{कुछ} \hi{लोगों} \hi{ने} \hi{अफ़वाह} \hi{से} \hi{भीड़} \hi{को} \hi{भड़काया।}
\end{quote}

\begin{quote}
\emph{kuch logoṁ ne afvāh se bhīṛ ko bhaṛkāyā.}
\end{quote}

\begin{quote}
A few people used the rumour to stir up the crowd.
\end{quote}

Past, with \hi{ने} (he or she …): \hi{उसने} \hi{लोगों} \hi{को} \hi{भड़काया}, \emph{usne logoṁ ko bhaṛkāyā}.

\begin{cousinweb}
\textbf{\hi{भड़काना}} (\emph{bhaṛkānā}) is the causative of \textbf{\hi{भड़कना}}, to flare up, from \emph{*bhaṭ} (CDIAL 9365): to make a fire, or a crowd, blaze.
\end{cousinweb}

\begin{grammarlens}[title={What you've built}]
One more word for telling what happened.
\end{grammarlens}

\subsection*{Your turn}
\begin{itemize}
\raggedright
  \item \emph{Say it:} \emph{bhaṛkānā}
  \item \emph{Say it:} \emph{bhaṛkānā}, once more
  \item \emph{Say it:} \emph{kuch logoṁ ne afvāh se bhīṛ ko bhaṛkāyā.} — A few people used the rumour to stir up the crowd.
  \item \emph{Recall:} say the Hindi for talk, then the Hindi for terror, then say \emph{bhaṛkānā} again
\end{itemize}

\begin{tcolorbox}[breakable,colback=teal!4,colframe=teal!35!black,title={Before you move on}]
What does \hi{भड़काना} mean? (\textquotedblleft{}To stir up; to incite\textquotedblright{}.) Say it once more.
\end{tcolorbox}
\section[\texorpdfstring{\hi{सनसनी} (sansanī)}{सनसनी (sansanī)}]{\hi{सनसनी} (sansanī) — a sensation; a thrill}

\label{lesson:HI-C314-sansani}



\begin{quote}
Before the new one: say the Hindi for terror, then the Hindi for to stir up.
\end{quote}

\subsection*{You'll want to know: \hi{सनसनी}}
\textbf{\hi{सनसनी}} — \emph{sansanī} — \textquotedblleft{}a sensation; a thrill\textquotedblright{}.

\begin{quote}
\hi{उस} \hi{ख़बर} \hi{से} \hi{पूरे} \hi{देश} \hi{में} \hi{सनसनी} \hi{फैल} \hi{गई।}
\end{quote}

\begin{quote}
\emph{us khabar se pūre deś meṁ sansanī phail gaī.}
\end{quote}

\begin{quote}
That news caused a sensation across the whole country.
\end{quote}

Feminine.

\begin{cousinweb}
\textbf{\hi{सनसनी}} (\emph{sansanī}) belongs with \emph{sansanānā}, to hum or tingle, which CDIAL 13901 sets under Sanskrit \emph{svana}, \textquotedblleft{}a noise\textquotedblright{}, noting that many of these words are imitative: a stir that runs through people like a buzz.
\end{cousinweb}

\begin{grammarlens}[title={What you've built}]
One more word for telling what happened.
\end{grammarlens}

\subsection*{Your turn}
\begin{itemize}
\raggedright
  \item \emph{Say it:} \emph{sansanī}
  \item \emph{Say it:} \emph{sansanī}, once more
  \item \emph{Say it:} \emph{us khabar se pūre deś meṁ sansanī phail gaī.} — That news caused a sensation across the whole country.
  \item \emph{Recall:} say the Hindi for terror, then the Hindi for to stir up, then say \emph{sansanī} again
\end{itemize}

\begin{tcolorbox}[breakable,colback=teal!4,colframe=teal!35!black,title={Before you move on}]
What does \hi{सनसनी} mean? (\textquotedblleft{}A sensation; a thrill\textquotedblright{}.) Say it once more.
\end{tcolorbox}
